% Source: https://artofproblemsolving.com/community/c3h1750568_problems_i_need_help_solving_asap
% Stable identity: aops-topic-1750568-q01
For $PC$ to cross side $AD$, $P$ must lie on line $AB$ beyond $A$. Use coordinates $A=(0,0)$, $B=(12,0)$, $C=(12,12)$, $D=(0,12)$, $P=(-4,0)$.

\begin{asy}
usepackage("amsmath");
size(6cm);
defaultpen(fontsize(10pt));
pair A=(0,0),B=(12,0),C=(12,12),D=(0,12),P=(-4,0),Q=(0,3);
draw(A--B--C--D--cycle,linewidth(0.9));
draw(P--A,linewidth(0.9));draw(P--C,red+linewidth(0.9));
dot(P);dot(Q);
label("$A$",A,S);label("$B$",B,SE);label("$C$",C,NE);label("$D$",D,NW);
label("$P$",P,SW);label("$Q$",Q,dir(140));
label("$4$",(P+A)/2,S);label("$12$",(A+B)/2,S);label("$12$",(B+C)/2,E);
label("$3$",(0,1.5),E);label("$9$",(0,7.5),E);
label("\mbox{}",(min(currentpicture,user=true)+max(currentpicture,user=true))/2,black); // keeps the text after the figure black when LaTeX typesets the labels inline
\end{asy}

Triangles $PAQ$ and $PBC$ are similar (both have a right angle, at $A$ and at $B$, and they share the angle at $P$), so
\[
\dfrac{AQ}{BC}=\dfrac{PA}{PB}=\dfrac{4}{16}\ \Longrightarrow\ AQ=3,\qquad DQ=9 .
\]
Then in right triangle $CDQ$, $CQ=\sqrt{CD^2+DQ^2}=\sqrt{144+81}=\boxed{15}$.
